\section{Discussion}
The evaluated daily workload--wellness panel did not demonstrate stable incremental forecasting value beyond a personalised readiness-history baseline. Across two within-team chronological evaluations and two bidirectional cross-team transportability evaluations, the full monitoring model did not lower the prespecified primary metric, MAE, after both comparators had been tuned with the same complete online deployment pipeline. Every player-cluster bootstrap interval for the primary MAE difference crossed zero. The direction of the findings was consistent across alternative residual-history windows, a reduced monitoring panel, and a wellness-item-missingness sensitivity analysis conditional on observed current- and next-day readiness. Because no smallest practically important MAE difference or decision threshold was specified, these findings do not establish formal equivalence, non-inferiority, or absence of all practical value.

The interpretation should remain narrow. The analysis does not show that training load, sleep, soreness, fatigue, stress, or mood are unimportant in athlete management. Such measures may facilitate communication, identify concerns not captured by a single readiness score, support clinical reasoning, or relate to subsequent performance and health outcomes. The present result addresses a specific conditional prediction task: whether the tested daily panel improved next-calendar-day self-reported readiness once the model already had access to current readiness and the same player's earlier prediction errors. In that setting, the added variables did not demonstrate a stable reduction in absolute error.

Several explanations are compatible with this result. First, current readiness may already integrate much of the recent workload and recovery information that players use to report short-term state. Second, the sequential residual adjustment may capture persistent player-specific reporting tendencies and unmeasured contextual influences. Third, individual variables may have useful signal for particular players or situations without providing enough stable, transportable information to improve pooled error across both teams. Finally, self-reported measures collected through the same monitoring system may share response and measurement characteristics, limiting the extent to which a broader questionnaire supplies independent forecast information.

The temporal design is a strength. Random record-level splits can yield optimistic estimates when observations from the same player and team context appear in both development and test samples \citep{Bergmeir2018,Roberts2017}. The present design separated source and target periods and, in two scenarios, source and target teams. The absence of stable incremental gain for \Pone{} under these conditions supports the use of strong personal-history comparators before additional monitoring variables are assumed to add short-horizon predictive value. The cross-team findings nevertheless require careful interpretation. They are a two-team transportability assessment under sequential target-team adaptation, not zero-shot external validation or evidence of universal external validity. In particular, the player-cluster bootstrap intervals quantify uncertainty across athletes within the evaluated target team, not uncertainty across a broad population of teams.

The findings also caution against equating a larger feature set with better short-horizon prediction. Session-RPE and wellness questionnaires can be easy to collect, but routine monitoring has opportunity costs. The appropriate implication is to align measurement burden with a defined decision problem, not to discontinue broader monitoring. When the decision is next-calendar-day readiness forecasting, future work should establish whether each additional measure improves a prespecified error target, calibration, decision utility, or a distinct clinical or performance endpoint.

Several limitations require emphasis. First, readiness was a self-reported 1--10 score rather than an objective performance, clinical, or injury endpoint. Second, the primary paired comparison required complete current-day wellness and observed next-day readiness, so it applies only to report-complete, outcome-observed days. The wellness-item-missingness sensitivity analysis added only 49 player-days and did not address dates on which the daily questionnaire, including current readiness, was not completed. Missingness may be associated with schedule, reporting behaviour, or athlete state; the present design cannot establish how model performance would translate to all calendar days. Third, the sample comprised two teams from one league using one monitoring protocol. Fourth, the analysis used subjective data only; objective external-load variables, match context, menstrual-cycle information, injury status, travel, and other contextual features could alter the incremental-value comparison. Fifth, the models were deliberately parsimonious linear/ridge implementations. Weekday and month were entered as numeric terms rather than as one-hot or cyclic functions, so nonlinear or cyclical calendar effects may not have been fully represented. Sixth, predictions were analysed on a continuous, unbounded scale; an operational system may require a prespecified range-handling rule and evaluation of its consequences. Seventh, the online residual adjustment assumes that earlier target-period readiness outcomes are available. It is appropriate for ongoing monitoring but does not represent a true new-player cold start. Finally, no smallest practically important MAE difference or decision threshold was prespecified. The analyses therefore do not establish formal equivalence, clinical utility, or the causal consequences of changing any monitoring variable \citep{Shmueli2010}.

Future studies should test the same conditional incremental-value question in independent women's football cohorts, include objective external-load and contextual features, explicitly model questionnaire non-completion, assess cold-start performance before residual histories accrue, compare alternative calendar encodings and prediction-boundary rules, prespecify decision-relevant performance thresholds, and evaluate decision-analytic utility. These studies should retain strict temporal alignment, participant-aware validation, calibration assessment, and explicit comparison with strong personal-history baselines \citep{Collins2024,Debray2023TRIPODCluster}.

\section{Conclusions}
For report-complete player-days in two elite women's football teams, the evaluated daily workload--wellness panel did not demonstrate stable incremental MAE improvement beyond a personalised readiness-history baseline with identical sequential residual updating. The result challenges the assumption that routinely adding more daily monitoring variables necessarily improves next-calendar-day self-reported readiness forecasting in this specific calendar-day deployment setting. It does not establish that workload or wellness measures are causally irrelevant, nor does it justify removing them from broader athlete-management practice.

\section*{Practical applications}
\begin{itemize}
\item Evaluate daily monitoring panels against a strong player-specific readiness-history comparator rather than a weak or randomly split benchmark.
\item Report the intended prediction time point, calendar-time validation, calibration, the scope of outcome and questionnaire availability, and the level at which uncertainty is quantified.
\item Treat additional monitoring variables as showing incremental predictive value only when they improve a prespecified decision-relevant outcome; they may still have separate clinical or communication value.
\item Do not infer from this study that teams should stop collecting workload or wellness information. The result is limited to conditional prediction of next-calendar-day self-reported readiness on report-complete days.
\end{itemize}
